% GENERATED FILE. Edit canonical lessons, then run npm run generate:books.
% canonical-source-hash: fnv1a64:f01c6a359b18cb15
% canonical-lessons: FR-C464-canicule, FR-C464-secheresse, FR-C464-humidite, FR-C464-moite, FR-C464-transpirer

\chapter{La canicule}
\label{ch:fr-la-canicule}
\hlchaptermodality{464}

\begin{chapteropening}
\textbf{By the end of this chapter:} \emph{I can talk about a heatwave: the heatwave, the drought, the humidity, clammy air and sweating.}

Complete the last lesson of chapter 464: I can talk about a heatwave: the heatwave, the drought, the humidity, clammy air and sweating.
\end{chapteropening}

\section[la canicule]{la canicule — the heatwave}

\label{lesson:FR-C464-canicule}



\begin{quote}
Before the new one: say the French for to thaw, then the French for to shiver.
\end{quote}

\subsection*{You'll want to know: la canicule}
\textbf{la canicule} — \textquotedblleft{}the heatwave\textquotedblright{}.

\begin{quote}
Pendant la canicule, les gens restent chez eux l'après-midi.
\end{quote}

\begin{quote}
\emph{During the heatwave people stay at home in the afternoon.}
\end{quote}

In Québec also \emph{la vague de chaleur}.

\begin{cousinweb}
\emph{canicule} is a learned borrowing from Latin \emph{canīcula} 'little dog', the name of the Dog Star, which rises in the hottest days of summer (EtymDB); it is built on \emph{canis} 'dog'.
\end{cousinweb}

\begin{grammarlens}[title={What you've built}]
One more word for telling, explaining and describing at B1.
\end{grammarlens}

\subsection*{Your turn}
\begin{itemize}
\raggedright
  \item \emph{Say it:} \emph{la canicule}
  \item \emph{Say it:} \emph{la canicule}, once more
  \item \emph{Say it:} a sentence of your own with \emph{la canicule}
  \item \emph{Recall:} say the French for to thaw, then the French for to shiver, then say \emph{la canicule} again
\end{itemize}

\begin{tcolorbox}[breakable,colback=teal!4,colframe=teal!35!black,title={Before you move on}]
What does la canicule mean? (\textquotedblleft{}The heatwave\textquotedblright{}.) Say it once more.
\end{tcolorbox}
\section[la sécheresse]{la sécheresse — the drought}

\label{lesson:FR-C464-secheresse}



\begin{quote}
Before the new one: say the French for to shiver, then the French for the heatwave.
\end{quote}

\subsection*{You'll want to know: la sécheresse}
\textbf{la sécheresse} — \textquotedblleft{}the drought\textquotedblright{}.

\begin{quote}
À cause de la sécheresse, les rivières sont presque vides.
\end{quote}

\begin{quote}
\emph{Because of the drought the rivers are almost empty.}
\end{quote}

Built on \emph{sec} 'dry'.

\begin{cousinweb}
\emph{sécheresse} is built on \emph{sec, sèche} 'dry' (EtymDB lists it), which is inherited from Latin \emph{siccus}.
\end{cousinweb}

\begin{grammarlens}[title={What you've built}]
One more word for telling, explaining and describing at B1.
\end{grammarlens}

\subsection*{Your turn}
\begin{itemize}
\raggedright
  \item \emph{Say it:} \emph{la sécheresse}
  \item \emph{Say it:} \emph{la sécheresse}, once more
  \item \emph{Say it:} a sentence of your own with \emph{la sécheresse}
  \item \emph{Recall:} say the French for to shiver, then the French for the heatwave, then say \emph{la sécheresse} again
\end{itemize}

\begin{tcolorbox}[breakable,colback=teal!4,colframe=teal!35!black,title={Before you move on}]
What does la sécheresse mean? (\textquotedblleft{}The drought\textquotedblright{}.) Say it once more.
\end{tcolorbox}
\section[l'humidité]{l'humidité — the humidity, the damp}

\label{lesson:FR-C464-humidite}



\begin{quote}
Before the new one: say the French for the heatwave, then the French for the drought.
\end{quote}

\subsection*{You'll want to know: l'humidité}
\textbf{l'humidité} — \textquotedblleft{}the humidity, the damp\textquotedblright{}.

\begin{quote}
Avec l'humidité, mes cheveux frisent.
\end{quote}

\begin{quote}
\emph{With the humidity my hair goes frizzy.}
\end{quote}

A feminine noun.

\begin{cousinweb}
\emph{humidité} is a learned borrowing from Latin \emph{hūmiditās} (EtymDB), from \emph{hūmidus} 'damp'.
\end{cousinweb}

\begin{grammarlens}[title={What you've built}]
One more word for telling, explaining and describing at B1.
\end{grammarlens}

\subsection*{Your turn}
\begin{itemize}
\raggedright
  \item \emph{Say it:} \emph{l'humidité}
  \item \emph{Say it:} \emph{l'humidité}, once more
  \item \emph{Say it:} a sentence of your own with \emph{l'humidité}
  \item \emph{Recall:} say the French for the heatwave, then the French for the drought, then say \emph{l'humidité} again
\end{itemize}

\begin{tcolorbox}[breakable,colback=teal!4,colframe=teal!35!black,title={Before you move on}]
What does l'humidité mean? (\textquotedblleft{}The humidity, the damp\textquotedblright{}.) Say it once more.
\end{tcolorbox}
\section[moite]{moite — clammy, sticky (with damp heat)}

\label{lesson:FR-C464-moite}



\begin{quote}
Before the new one: say the French for the drought, then the French for the humidity.
\end{quote}

\subsection*{You'll want to know: moite}
\textbf{moite} — \textquotedblleft{}clammy, sticky (with damp heat)\textquotedblright{}.

\begin{quote}
Dans le métro, l'air est chaud et moite.
\end{quote}

\begin{quote}
\emph{In the metro the air is hot and clammy.}
\end{quote}

\emph{Les mains moites} — sweaty hands. Same form for both genders.

\begin{cousinweb}
\emph{moite} comes from Old French \emph{moiste} (EtymDB), which is usually traced to Latin \emph{mūcidus} 'mouldy, slimy' (Lewis \& Short); the details are not settled.
\end{cousinweb}

\begin{grammarlens}[title={What you've built}]
One more word for telling, explaining and describing at B1.
\end{grammarlens}

\subsection*{Your turn}
\begin{itemize}
\raggedright
  \item \emph{Say it:} \emph{moite}
  \item \emph{Say it:} \emph{moite}, once more
  \item \emph{Say it:} a sentence of your own with \emph{moite}
  \item \emph{Recall:} say the French for the drought, then the French for the humidity, then say \emph{moite} again
\end{itemize}

\begin{tcolorbox}[breakable,colback=teal!4,colframe=teal!35!black,title={Before you move on}]
What does moite mean? (\textquotedblleft{}Clammy, sticky (with damp heat)\textquotedblright{}.) Say it once more.
\end{tcolorbox}
\section[transpirer]{transpirer — to sweat, to perspire}

\label{lesson:FR-C464-transpirer}



\begin{quote}
Before the new one: say the French for the humidity, then the French for clammy.
\end{quote}

\subsection*{You'll want to know: transpirer}
\textbf{transpirer} — \textquotedblleft{}to sweat, to perspire\textquotedblright{}.

\begin{quote}
Il transpire beaucoup quand il fait chaud.
\end{quote}

\begin{quote}
\emph{He sweats a lot when it is hot.}
\end{quote}

A little more polite than \emph{suer}.

\begin{cousinweb}
\emph{transpirer} is a learned borrowing from Medieval Latin \emph{transpīrāre} 'to breathe through' (EtymDB), from \emph{trāns} and \emph{spīrāre} 'to breathe'.
\end{cousinweb}

\begin{grammarlens}[title={What you've built}]
One more word for telling, explaining and describing at B1.
\end{grammarlens}

\subsection*{Your turn}
\begin{itemize}
\raggedright
  \item \emph{Say it:} \emph{transpirer}
  \item \emph{Say it:} \emph{transpirer}, once more
  \item \emph{Say it:} a sentence of your own with \emph{transpirer}
  \item \emph{Recall:} say the French for the humidity, then the French for clammy, then say \emph{transpirer} again
\end{itemize}

\begin{tcolorbox}[breakable,colback=teal!4,colframe=teal!35!black,title={Before you move on}]
What does transpirer mean? (\textquotedblleft{}To sweat, to perspire\textquotedblright{}.) Say it once more.
\end{tcolorbox}
